% Part: computability
% Chapter: recursive-functions
% Section: composition

\documentclass[../../../include/open-logic-section]{subfiles}

\begin{document}

\olfileid{cmp}{rec}{com}
\olsection{复合}

如果 $f$ 和 $g$ 是两个自然数的一元函数，我们可以将它们复合：$h(x) = f(g(x))$。这样，新函数~$h(x)$ 就是由函数 $f$ 和~$g$ 通过\emph{复合}定义的。我们希望将这个概念推广到多于一个自变元的函数。

推广的一种方式如下：假设 $f$ 是一个 $k$ 元函数，而 $g_0$，……，$g_{k-1}$ 是 $k$ 个 $n$ 元函数。那么我们可以定义一个新的 $n$ 元函数 $h$ 如下：
\[
h(x_0, \dots, x_{n-1}) =
f(g_0(x_0, \dots, x_{n-1}), \dots, g_{k-1}(x_0, \dots, x_{n-1}))
\]
如果 $f$ 和所有 $g_i$ 都是可计算的，那么 $h$ 也是可计算的：要计算 $h(x_0,
\dots, x_{n-1})$，首先对每个 $i = 0$，……，$k-1$ 计算出 $y_i = g_i(x_0, \dots,
x_{n-1})$ 的值。然后将这些值代入 $f$ 来计算 $h(x_0, \dots, x_{k-1}) = f(y_0, \dots, y_{k-1})$。

这看起来可能过于严格地刻画了我们利用已有函数计算新函数时的做法。一方面，有时我们并不使用一个函数的所有自变元，就像我们在~$\Add$ 的原始递归定义中定义 $g(x, y, z) = \Succ(z)$ 以供使用时那样。假设我们允许使用以下函数：
\[
\Proj{n}{i}(x_0, \dots, x_{n-1}) = x_i
\]
函数~$\Proj{k}{i}$ 称为\emph{投影}函数：$\Proj{n}{i}$ 是一个 $n$ 元函数。那么 $g$ 就可以定义为
\[
g(x, y, z) = \Succ(\Proj{3}{2}(x, y, z)).
\]
这里一元函数 $\Succ$ 扮演 $f$ 的角色，所以 $k=1$。而三元函数 $\Proj{3}{2}$ 扮演 $g_0$ 的角色。结果是一个三元函数，它返回第三个自变元的后继。

投影函数也允许我们通过重排或等同自变元来定义新函数。例如，函数 $h(x)
= \Add(x, x)$ 可以定义为
\[
h(x_0) = \Add(\Proj{1}{0}(x_0),\Proj{1}{0}(x_0)).
\]
这里 $k=2$，$n=1$，$\Add$ 扮演 $f(y_0,y_1)$ 的角色，而~$\Proj{1}{0}(x_0)$（即一元投影函数，也称为恒等函数）同时扮演了 $g_0(x_0)$ 和 $g_1(x_0)$ 的角色。

如果我们已有一个函数 $f(y_0, y_1)$，我们可以通过下式定义函数 $h(x_0, x_1) = f(x_1, x_0)$：
\[
h(x_0, x_1) = f(\Proj{2}{1}(x_0, x_1),\Proj{2}{0}(x_0, x_1)).
\]
这里 $k=2$，$n = 2$，而 $g_0$ 和 $g_1$ 的角色分别由 $\Proj{2}{1}$ 和~$\Proj{2}{0}$ 扮演。

你可能还会担心，$g_0$，……，$g_{k-1}$ 都被要求具有相同的元数~$n$。（记住，函数的\emph{元数}是其自变元的个数；一个 $n$ 元函数的元数是~$n$。）但引入投影函数提供了所需的灵活性。例如，假设 $f$ 和~$g$ 是三元函数，而 $h$~是由下式定义的二元函数：
\[
h(x,y) = f(x,g(x,x,y),y).
\]
利用投影函数，$h$ 的定义可以重写为：
\[
h(x,y) = f(\Proj{2}{0}(x,y), g(\Proj{2}{0}(x,y), \Proj{2}{0}(x,y),
\Proj{2}{1}(x,y)), \Proj{2}{1}(x,y)).
\]
那么 $h$ 是 $f$ 与 $\Proj{2}{0}$、$l$ 和 $\Proj{2}{1}$ 的复合，其中
\[
l(x,y) = g(\Proj{2}{0}(x,y),\Proj{2}{0}(x,y),\Proj{2}{1}(x,y)),
\]
也就是说，$l$ 是 $g$ 与 $\Proj{2}{0}$、$\Proj{2}{0}$ 和~$\Proj{2}{1}$ 的复合。

\end{document}